\section{An order bound for multiplicative-shift systems}
\label{sec:order}

A bound on the order of a nonzero scalar series produced by the formula
system will make a finite truncation possible.  We prove such a bound for
a general rational system under a multiplicative shift.  We adapt the
derivative-row and minor comparison underlying Moura's multiplicity
estimate~\cite{Moura2003}, in the formal presentation of
\cite[Section~4, Theorem~4.2]{OAInc}, to successive shifts.
The distinct-order basis and leading-term calculation parallel the
Wronskian argument of Bostan and Dumas~\cite[Lemmas~1--3]{BostanDumas2010}.
Here the leading determinant uses distinct powers of the shift parameter,
and remains nonzero in every characteristic.

For a nonzero series $h\in K[[z]]$, write $\ord_z h$ for its least exponent
with nonzero coefficient, and put $\ord_z 0=+\infty$.

\begin{theorem}[Order bound]\label{thm:order}
Let $K$ be a field, let $\tau\in K^\times$ have infinite multiplicative
order, and let $\sigma$ be the automorphism of $K((z))$ that fixes $K$ and
sends $z$ to $\tau z$.  Let $m\geq 1$ and $\nu\geq 0$ be integers.
Suppose that $H\in\GL_m(K[[z]])$ satisfies
\begin{equation}\label{eq:order-system}
  \sigma(H)=A H,\qquad A=\frac{P}{Q},
\end{equation}
where $P\in\Mat_m(K[z])$, $Q\in K[z]$, $Q(0)\neq 0$, and $Q$ and every
entry of $P$ have degree at most $\nu$.  For every constant row
$u\in K^{1\times m}$ and constant column $v\in K^{m\times 1}$,
\begin{equation}\label{eq:order-bound}
  uHv\neq 0
  \quad\Longrightarrow\quad
  \ord_z(uHv)\leq \nu\binom{m}{2}.
\end{equation}
\end{theorem}

\begin{proof}
Set $g=uH$ and suppose that $y=gv\neq 0$.  Let $V\subseteq K[[z]]$ be the
$K$-linear span of the $m$ components of $g$, and put $k=\dim_K V$.
Then $1\leq k\leq m$ and $y\in V$.  We first express the orders occurring
in $V$ through a determinant, and then bound that determinant using the
rational system.

There is a basis $w_1,\ldots,w_k$ of $V$ whose orders satisfy
\[
  0\leq c_1<c_2<\cdots<c_k,
  \qquad c_i=\ord_z w_i.
\]
To construct it, choose a nonzero element of smallest order in $V$.
The coefficient functional at that order is nonzero, so its kernel has
codimension one and every nonzero element of the kernel has larger
order.  Repeat inside this kernel, and continue until the dimension is
zero.  In any nonzero linear combination of the resulting basis, the
first basis element with a nonzero coefficient supplies the unique
lowest-order term.  In particular,
\begin{equation}\label{eq:order-largest-basis-order}
  \ord_z y\leq c_k.
\end{equation}

Consider the determinant of successive shifts, or Casoratian,
\[
  C(w_1,\ldots,w_k)
  =\det\bigl(\sigma^j(w_i)\bigr)_{
      \substack{0\leq j<k\\1\leq i\leq k}}.
\]
Write the leading term of $w_i$ as $\lambda_i z^{c_i}$, with
$\lambda_i\in K^\times$.  After factoring $z^{c_i}$ from column $i$, the
constant term of the determinant that remains is
\[
  \left(\prod_{i=1}^k\lambda_i\right)
  \det\bigl((\tau^{c_i})^j\bigr)_{
      \substack{0\leq j<k\\1\leq i\leq k}}
  =\left(\prod_{i=1}^k\lambda_i\right)
    \prod_{1\leq i<h\leq k}(\tau^{c_h}-\tau^{c_i}).
\]
This expression is nonzero because $\tau$ has infinite multiplicative
order.  Therefore
\begin{equation}\label{eq:casoratian-order}
  \ord_z C(w_1,\ldots,w_k)=\sum_{i=1}^k c_i.
\end{equation}
Every other $K$-basis of $V$ gives a Casoratian of the same order: its
matrix of successive shifts is obtained by an invertible change of
columns over $K$, which is fixed by $\sigma$.

To bound this order, we compare the shifted components of $g$ with
rational rows obtained from the system.
Define rational rows
\begin{equation}\label{eq:order-rational-rows}
  R_0=u,\qquad R_{j+1}=\sigma(R_j)A\quad(j\geq 0).
\end{equation}
Equation~\eqref{eq:order-system} gives $\sigma^j(g)=R_jH$ by induction.
Each row has a presentation
\begin{equation}\label{eq:order-row-denominators}
  R_j=\frac{N_j}{D_j},\qquad
  D_j(z)=\prod_{e=0}^{j-1}Q(\tau^e z),
\end{equation}
where $N_j$ is a polynomial row whose entries have degree at most $j\nu$;
the empty product is $D_0=1$.  Indeed, take $N_0=u$ and then
$N_{j+1}=\sigma(N_j)P$.
Substitution by $\tau z$ preserves polynomial degree, which proves the
degree bound.  Moreover, $D_j(0)=Q(0)^j\neq 0$, so every entry of every
$R_j$ belongs to $K[[z]]$.

Let $M$ be the $k\times m$ matrix with rows $R_0,\ldots,R_{k-1}$.
Its product $MH$ has rows $g,\sigma(g),\ldots,\sigma^{k-1}(g)$.
For a $k\times m$ matrix $X$ and a $k$-element subset
$J\subseteq\{1,\ldots,m\}$, let $X_J$ denote the square matrix formed
by the columns indexed by $J$, in increasing order.  A determinant
$\det((MH)_J)$ is the Casoratian of the corresponding components of
$g$.  It vanishes when those components are $K$-linearly dependent.
Otherwise they form a basis of $V$, and
\eqref{eq:casoratian-order} shows that its order is $\sum_i c_i$.
Some $k$ components form a basis of $V$, so at least one of these
determinants is nonzero.  Consequently,
\begin{equation}\label{eq:order-casoratian-minors}
  \min_J\ord_z\det((MH)_J)=\sum_{i=1}^k c_i.
\end{equation}

Right multiplication by $H$ preserves the minimum order of the
$k\times k$ minors.  To see this directly, Cauchy--Binet expresses every
minor of $MH$ as a $K[[z]]$-linear combination of minors of $M$.
The minimum order therefore cannot decrease.  Applying the same
argument to $M=(MH)H^{-1}$ proves the reverse inequality, since
$H^{-1}$ also has power series entries.  Thus
\begin{equation}\label{eq:order-minor-invariance}
  \min_J\ord_z\det(M_J)
  =\min_J\ord_z\det((MH)_J).
\end{equation}
In particular, $M$ has a nonzero $k\times k$ minor.

By~\eqref{eq:order-row-denominators}, any such nonzero minor has
denominator $\prod_{j=0}^{k-1}D_j$, a unit at zero, and a nonzero
polynomial numerator of degree at most
\[
  \sum_{j=0}^{k-1}j\nu=\nu\binom{k}{2}.
\]
Its order is at most this degree.  Combining
\eqref{eq:order-largest-basis-order},
\eqref{eq:order-casoratian-minors}, and
\eqref{eq:order-minor-invariance} gives
\[
  \ord_z y
  \leq c_k
  \leq\sum_{i=1}^k c_i
  \leq\nu\binom{k}{2}
  \leq\nu\binom{m}{2},
\]
as claimed.
\end{proof}

The characteristic plays no role in the Vandermonde calculation:
distinct orders $c_i$ yield distinct powers $\tau^{c_i}$.  In the
application, $K=F(t)$ and $\tau=t$, so this condition holds for every
coefficient field $F$.  The other essential local condition is that
$H$ be invertible over $K[[z]]$, as required to use both $H$ and
$H^{-1}$ in~\eqref{eq:order-minor-invariance}.
